\documentclass[12pt,a4paper,twoside]{article}

\usepackage{parskip}
\usepackage{ebgaramond}
\usepackage{csquotes}
\usepackage{booktabs}


\title{Geomerge}
\author{Vincent Liu}

\begin{document}

\maketitle

\section{Introduction}

The purpose of this doc is to expand on the operation log part of the original
Geolog prospectus with more details on what/how we might want to store in the
operation log. Quoting from the prospectus:

\begin{displayquote}
An operation log which stores a (not-necessarily linear) history of all operations
performed on a database instance.
\end{displayquote}

\begin{displayquote}
  The ``operation log'' part is intended to be very similar to current data structures
  inside Automerge.
\end{displayquote}

And we should do so by looking at a concrete example of a problem that can be
expressed in Geolog. This has the danger of overfitting our design to it, so
this doc would evolve as we discover more. But my hope is that this can inform
discussions both across and within different teams working on Geolog.

\section{Motivating example}

Let us consider this graph example from geolog-zeta:

\begin{verbatim}
  theory Graph {
    V : Sort;
    E : Sort;
    src : E -> V;
    tgt : E -> V;
    reachable : [from: V, to: V] -> Prop;
  
    ax/edge : forall e : E. |- [from: e src, to: e tgt] reachable;
    ax/trans : forall x,y,z : V.
      [from: x, to: y] reachable, [from: y, to: z] reachable
      |- [from: x, to: z] reachable;
  }
  
  
  instance G : Graph = chase {
    a, b, c : V;
    e1, e2 : E;
    e1 src = a; e1 tgt = b;
    e2 src = b; e2 tgt = c;
  } 
\end{verbatim}

Here we have a theory \verb|Graph| which defines data of different sorts (\verb|V| and
\verb|E|), functions (\verb|src| and \verb|tgt|) and relations (\verb|reachable|). 
And two axioms written as sequents (i.e.\ geometric theory) that defines what 
must hold for the data, functions and relations.

Then we have a concrete instance \verb|G| that gives the sorts and functions, but not
relations. If we did not include \verb|chase| in our definition, this will be rejected
by the elaborator because \verb|ax/edge| is not satisfied.

However, if we do include \verb|chase|, then the chase algorithm will compute the
fixpoint for us:

\begin{verbatim}
  instance G : Graph = {
  // V (3):
  a : V;
  b : V;
  c : V;
  // E (2):
  e1 : E;
  e2 : E;
  // src:
  e1 src = a;
  e2 src = b;
  // tgt:
  e1 tgt = b;
  e2 tgt = c;
  // reachable (3 tuples):
  [from: a, to: b] reachable;
  [from: b, to: c] reachable;
  [from: a, to: c] reachable;
} 
\end{verbatim}

Now what exactly do we want to store in our operation log? We can encode the theory
and instances. But I imagine the instance is where there will be lots of data
and ``operations'', such as adding and retracting elements from it, etc.

\section{Operation log content}

To understand what we want to store in our operation log, we need to first define
what an operation is. Looking at the geolog-zeta, there are three supported operations:
`add`, `retract` and `assert`, corresponding to adding an element, removing an
element, and adding a new relation. These are viable candidates to be defined
as our operations. And the idea is that, because our log is monotonic (need to
be careful with retraction), we can achieve incremental evaluation, by, say,
applying Datalog style semi-naive evaluation strategy where only the new elements
added need to be evaluated against the existing data.

By utilising Automerge's opset design, this would give us the ability to build
concurrency into Geolog, where multiple agents/people/threads can perform these
operations on an instance concurrently, and exchange new operations when needed
to (e.g. they wish to merge). And thanks to monotonicity, when receiving the new 
operations, the query engine only can apply incremental evaluation.

As an example, the simplest case one might imagine is adding another element \verb|d:V|
to instance \verb|G| and an edge \verb|e3:E| with \verb|e3 src: c| and \verb|e3 tgt: d|. And if we
want to query the \verb|reachable| relation, we just need to run the axioms on our
existing data against the newly added data, and repeat until we reach the fixpoint.
The new \verb|reachable(c, d), reachable(a, d), reachable(b, d)| should be added
to the instance data as an operation \verb|assert| as well (?).

Now let's consider what happens when two agents merge their instances. If we disallow
retraction, this will be trivial, since we can simply merge the new data, and run
our chase algorithm on the new data until fix point. This behaviour would be similar
to a grow-only set CRDT (G-Set).
 
\section{Operation log format}

The format of the operation log will be similar to what already exists in Automerge.
Automerge already supports three types of objects: Map, List and Text. The datamodel
used by Geolog behaves like a set, which is more or less a boolean valued Map.
So we could easily use Automerge's opset to do the job, along with columnar compression,
and other nice features.

As an illustration, this is an example of what the opset might look like:


\begin{table}
\footnotesize
\begin{tabular}{cccccccccccccc}\toprule
  \multicolumn{2}{c}{objectId} & \multicolumn{3}{c}{key} & \multicolumn{2}{c}{opId} & action & insert & \multicolumn{2}{c}{value} & \multicolumn{3}{c}{succs} \\
  \cmidrule(lr){1-2}\cmidrule(lr){3-5}\cmidrule(lr){6-7}\cmidrule(lr){10-11}\cmidrule(lr){12-14}
  int & actor & str       & int & actor & int   & actor &        &       & type    & raw    & num   & int & actor \\\midrule
  -   & -     & instance  & -   & -     & \(1\) & A     & mkInst & false & -       & -      & \(0\) &     & \\
  1   & A     & d         & -   & -     & 2     & B     & add    & false & V       & -      & \(0\) & -   & - \\
  1   & A     & e         & -   & -     & 3     & A     & add    & false & E       & -      & \(0\) & -   & - \\
  1   & A     & reachable & -   & -     & 4     & A     & assert & false & RelArgs & (a, b) & \(0\) & -   & - \\\bottomrule
\end{tabular}
\end{table}

\begin{itemize}
  \item objectId would be the instance/theory id. For our case, we should consider each instance/theory 
  as an object
  \item key is a 3-tuple, where the str would be the name of element being added, 
  e.g.\ \verb|d| would be the name of the element added\@. (int,actor) here is
  not relevant.
  \item opId is the unique operation identifier. Here we are using the clock at
  each node.
  \item action would consist of the three instance mutation opertaion: add, retract and assert
  (how to add function definitions?)
  \item insert is used for operations on a list, not relevant here
  \item value is a 2-tuple, where the first element specifies the type of the element,
  if this was adding an element, that would be its sort. If it's asserting a relation,
  then it would be \verb|RelArgs|. Although I am not sure if we should store all the payload
  of the relation in this single column. We should probably store a pointer (foreign key)
  of some kind that points to the actual storage of the arguments of the tuple.
  \item succ is for dependency tracking, and is only needed when two actors make
  operations on the same element.
\end{itemize}

This is a small change to what the is already supported in Automerge's opset,
essentially the table format remains the same, but some of the fields have been
modified slightly. I want to use this as a starting point, and if we have more
real world use cases, we can consider redesigning as necessary. 


\section{Problems to think about}

\begin{itemize}
  \item When the user declares an instance with some initial elements in it, do we
  just store them as if these are added one by one?
  \item Is there anything that the opset needs to do to help with the query engine?
  What would their interaction be like?
\end{itemize}


% Regarding the interaction between Geomerge and Geodeed (query engine). I believe 
% there is not much, Geomerge is concerned about tracking dependencies of the operation
% history, and provide the right dependency (or the notion of time) to the upper layer,
% while Geodeed would be responsible for tracking the edb and idb and perform the
% correct incremental evaluation.


\bibliographystyle{plainnat}
\bibliography{geolog.bib}

\end{document}
